\section{Diagnostic Design and Method Space}

SID papers now optimize different artifact properties: learned item indexing
and TIGER-style generative retrieval established the interface~\cite{hua2023indexids,rajput2023tiger},
while later work stresses collaborative alignment, end-to-end or
differentiable tokenization, representation-source sensitivity,
code-assignment diversity, collision qualification, variable-length
interfaces, drift, and deployment surfaces~\cite{zheng2023lcrec,zhu2024cost,wang2024letter,liu2024elit,fu2026diger,qiao2026textasvision,hu2026quasid,pan2026adasid,cheng2026capsid,baikalov2026staleness,ju2026snapchatsid}.
The literature space is broader than the artifacts that v0 can honestly
execute. Table~\ref{tab:method-coverage} therefore reports only rows that pass
the validator or are explicitly scoped as controls, stressors, or backlog. The
method taxonomy guides future adapters; it is not presented as reproduced
method coverage.

\begin{table*}[t]
\caption{Evidence map for the v0 artifact. Literature method-space coverage is
discussed in the text; this table reports only artifacts that pass the v0
validator, controlled rows used for diagnostic calibration, and backlog rows
that do not support empirical claims.}
\label{tab:method-coverage}
\scriptsize
\setlength{\tabcolsep}{2.0pt}
\begin{tabular}{@{}>{\raggedright\arraybackslash}p{0.19\textwidth}>{\raggedright\arraybackslash}p{0.10\textwidth}>{\raggedright\arraybackslash}p{0.14\textwidth}>{\raggedright\arraybackslash}p{0.11\textwidth}>{\centering\arraybackslash}p{0.05\textwidth}>{\raggedright\arraybackslash}p{0.10\textwidth}>{\raggedright\arraybackslash}p{0.15\textwidth}@{}}
\toprule
Row & Cluster & Status & Items & Seeds & Diagnostics & Caveat \\
\midrule
GRID/RQ-KMeans feature-text & A & real export & 23,742 / 50,000 & 3 / 1 & D1--D5a & processed feature text; not raw-text TIGER/GRID \\
ReSID/GAOQ bounded & B1/B2 & real export & 23,742 & 1 & D1--D5a & one-epoch bounded Musical mapping \\
Category-prefix & control & built-in & 23,742 & n/a & D1--D5a & deterministic sanity row \\
Mod-collision hash & control & built-in & 23,742 & n/a & D1--D5a & collision lower-bound control \\
Popularity-balanced & control & built-in & 23,742 & n/a & D1--D5a & popularity calibration row \\
Qualified collision & stressor & built-in & synthetic & n/a & D2b,D3 & not named-method evidence \\
Capacity budget & stressor & built-in & synthetic & n/a & D1,D2,D4,D5a & not named-method evidence \\
Variable depth & stressor & built-in & synthetic & n/a & D5a & optional boundary evidence \\
15+ named SID lines & backlog & not run & -- & -- & -- & adapters require official artifacts \\
\bottomrule
\end{tabular}
\end{table*}

\paragraph{D1--D4.}
D1 reports per-level usage, prefix counts, imbalance summaries, and dead or
underused code indicators. D2 measures duplicate full \sid{}s and prefix
collisions; it is a collision profile, not a causal harm estimate. A bounded
interaction-qualified controller is included because collision-aware work
argues that collisions are not uniformly harmful~\cite{hu2026quasid}. D3 asks
whether prefix neighborhoods recover item co-occurrence neighbors, keeping
metadata purity as auxiliary context rather than as collaborative quality. D4
splits items by popularity and measures whether full-code capacity and prefix
structure are allocated differently across head, mid, and tail items.

\paragraph{D5a--D7.}
D5a reports structural cost proxies: \sid length, prefix fan-out, duplicate
codes, and trie-like expansion pressure. These are not serving latencies; long
or variable SID methods make this distinction important~\cite{xia2026acerec,cheng2026capsid}.
D6 computes churn between tokenizer refreshes for drift-aware settings~\cite{feng2026dact,baikalov2026staleness}.
D7 is an interface hook for generator behavior, such as invalid paths,
next-token entropy, and duplicate generated candidates; it requires
\texttt{generator\_outputs} and is not claimed by the current evidence.
